% Fuente: Auxiliar 3, sección «Poliedro».
Considera el poliedro en forma estándar \( P = \{x \mid Ax = b, \, x \geq 0\} \). Suponga que la matriz \( A \) tiene dimensiones \( m \times n \) y que sus filas son linealmente independientes. Para cada una de las siguientes afirmaciones, indica si es verdadera o falsa. Si es verdadera, explique por qué lo es (puede ser útil un gráfico); de lo contrario, proporciona un contraejemplo.

\begin{enumerate}
    \item[(a)] Si \( n = m + 1 \), entonces \( P \) tiene como máximo dos soluciones básicas factibles.
    \item[(b)] El conjunto de todas las soluciones óptimas es acotado.
    \item[(c)] En toda solución óptima, no puede haber más de \( m \) variables positivas.
    \item[(d)] Si hay más de una solución óptima, entonces hay incontablemente muchas soluciones óptimas.
    \item[(e)] Si existen varias soluciones óptimas, entonces existen al menos dos soluciones básicas factibles que son óptimas.

\end{enumerate}

